\documentclass{article}
\usepackage{amsthm, amssymb, latexsym, amsmath, amsbsy, graphicx, epsfig}
\usepackage{ mathrsfs }

\title{Writing Assignment \#11: Functions, Inverses, and Modular Arithmetic}
\author{Robert T Kohn}
\date{April 2025}

\begin{document}

\maketitle
{\huge{Section 4.2}}
\section*{Question 1}
\subsection*{Part D}

The sine function is not bijective as it is not one to one. However, over a closed domain from $\frac{-\pi}{2},\frac{\pi}{2}$ it does work because there are no repeated outputs. The inverse would be the arcsin of x when restricted.

\subsection*{Part F}

This function can be simplified to 1-$\frac{1}{x}$ in order to find the inverse. 

$y-1=-\frac{1}{x} \rightarrow 1-y=\frac{1}{x} \rightarrow x=\frac{1}{1-y}$

This is a function, for all x not equal to 0, as it is bijective. 

\section*{Question 12}

Claim: $\exists f \ni Dom(D) \wedge D \subseteq A \wedge f=g|_D \wedge \forall x, i(x)=x \Rightarrow f= g$ \textopenbullet  $i$

Then, $(\forall x \in D) \ni f(x)=g(x)$

Thus, $(\forall x \in i(x)=x) \ni g(i(x))=g(x)$

So, f(x)=g(i(x))

Hence, f(x)=(g \textopenbullet  i)(x)

Then, $f=g$\textopenbullet$i|_D$,

Thus, f=g \textopenbullet i

$\therefore \exists f \ni Dom(D) \wedge D \subseteq A \wedge f=g|_D \wedge \forall x, i(x)=x \Rightarrow f= g$ \textopenbullet  $i$\

\# this feels really messy but I didn't really know what else to do here, I think this works?

\section*{Question 15}
\subsection*{Part A}
Claim: $f: A \rightarrow B \wedge g: C \rightarrow D \wedge f \times g = \{(a,c),(b,d): (a,b) \in f \wedge (c,d) \in g \} \Rightarrow f \times g \ is \ a \ mapping \ from \ A \times C \ to \ B \times D$

Suppose $(a,c) \in A \times C$:

Then, $(\exists b \in B)(\exists d \in D) \ni (a,b) \in f \wedge (c,d) \in g$

Thus, ((a,c),(b,d)) $\in f \times g$

So, $(\forall(a,c) \in A \times C)(\exists(b,d)\in B \times D) \ni  ((a,c),(b,d)) \in f \times g$

$\therefore f: A \rightarrow B \wedge g: C \rightarrow D \wedge f \times g = \{(a,c),(b,d): (a,b) \in f \wedge (c,d) \in g \} \Rightarrow f \times g \ is \ a \ mapping \ from \ A \times C \ to \ B \times D$

\subsection*{Part B}

(f(a),g(c))

\section*{Question 19}


\subsection*{Part B}

This proof gets a grade of A.

\subsection*{Part C}

This proof gets a grade of F because it assumes that g \textopenbullet f = f \textopenbullet g which is not always true.

\

{\huge{Section 4.3}}
\section*{Question 3}
\subsection*{Part B}

This function is surjective as each y cooresponds to an x but it is not injective because multiple integers can map to $\mathbb{Z}_5$

\subsection*{Part C}

Same type of logic applies to this one. This is surjective as each output has an input but each output has duplicate inputs.

\section*{Question 9}
\subsection*{Part B}

$A=\{1\}, B=\{2,3\}, C=\{4,5\}$

f(1)=2, g(2)=4, g(3)=5

g \textopenbullet f(1)=g(2)=4

\subsection*{Part D}

$A=\{1,2\}, B=\{3,4,5\}, C=\{6\}$

f(1)=3, f(2)=4, g(3)=6, g(4)=6, g(5)=6

g \textopenbullet f(1)=g(3)=6, g \textopenbullet f(2)=g(4)=6

\section*{Question 12}
\subsection*{Part B}

Claim: f: $\mathbb{Z}_4 \rightarrow \mathbb{Z}_2$ given by $f(\overline{x})=[3x]$ is onto but not one to one.

Evaluating:

$f(\overline{0})=[0], f(\overline{1})=[3]=[1], f(\overline{2})=[6]=[0], f(\overline{3})=[9]=[1]$

Onto:

Then, both $[0] \wedge [1] \in \mathbb{Z}_2$ are achieved

Thus, every element of the codomain has a preimage

So,  f is onto

Not one to one:

Then, $f(\overline{0})=f(\overline{2}=[0]$

So, f maps distinct elements to the same value

Hence, f is not one to one.

$\therefore f: \mathbb{Z}_4 \rightarrow \mathbb{Z}_2$ given by $f(\overline{x})=[3x]$ is onto but not one to one.

\subsection*{Part D}
Claim: f: $\mathbb{Z}_4 \rightarrow \mathbb{Z}_4$ given by $f(\overline{x})=[2x]$ is neither onto nor not one to one.

Evaluating:
$f(\overline{0})=\overline{0}, f(\overline{1})=\overline{2}, f(\overline{2})=\overline{0} f(\overline{3})=\overline{2}$

Not one to one:

Then, $f(\overline{0})=f(\overline{2})=\overline{0}$

Thus, distinct inputs have the same output

So, f isn't one to one

Not onto:

Then, the image of f is $\{\overline{0},\overline{2}\}$ 

So, $\overline{1} \vee \overline{3}$ are not in the image

Thus, f isn't onto

$\therefore f: \mathbb{Z}_4 \rightarrow \mathbb{Z}_4$ given by $f(\overline{x})=[2x]$ is neither onto nor not one to one.

\section*{Question 15}
\subsection*{Part C}

This proof gets a grade of C as it states a correct idea and process but fails to prove surjectivity as it doesn't refer to arbitrary elements in C.

\subsection*{Part G}

This proves gets a grade of C as it incorrectly uses injectivity.

\end{document}
